\begin{tabularx}{\textwidth}{L{31mm}L{48mm}X}
\toprule
\textbf{Rész} & \textbf{Párhuzamosítás} & \textbf{Korlát / megjegyzés} \\
\midrule
Rekurzív felosztás & A bal és jobb résztömb rendezése egymástól függetlenül indítható. & Jól illeszkedik az oszd meg és uralkodj elvhez, de a szintek végén összefésülés szükséges. \\
Összefésülés & A másolás és bizonyos összehasonlítási részfeladatok párhuzamosíthatók. & A hagyományos kétmutatós összefésülés alapvetően lineáris és függőségeket tartalmaz. \\
Batcher-féle páros-páratlan összefésülés & Az elemek paritás szerinti bontásával rekurzív, hálózatszerű összefésülés építhető. & Jobban párhuzamosítható, de több összehasonlítást és bonyolultabb szerkezetet eredményezhet. \\
Gyakorlati küszöb & Kis résztömböket gyakran sorosan rendezünk. & A túl sok apró párhuzamos feladat adminisztrációs költsége ronthatja a gyorsítást. \\
\bottomrule
\end{tabularx}
